\section{Introduction}
\label{sec:introduction}

The wings of an aircraft are designed to withstand a wide range of operating conditions, but a small subset of events ends up driving the structural sizing. Among these, atmospheric gusts play a central role~\cite{hoblit1988gust,wright2007aeroelasticity}. A vertical gust hitting the wing produces a sudden change in the angle of attack, which generates a peak in the lift and a transient oscillation of the structure. The amplitude of this peak, together with the bending moment that follows, contributes to the design loads that the wing has to carry without failure, and the certification specifications prescribe the discrete gust and continuous turbulence conditions that the airframe must survive~\cite{EASA_CS25}. A wing sized only on this basis tends to be heavy, because the structure is dimensioned against a load that occurs for a very small fraction of the flight envelope.


\begin{figure}[t]
    \centering % Centra l'immagine nella pagina
    \includegraphics[width=\textwidth,height=1\textheight,keepaspectratio]{Thesis_article_format/Chapters/Chapter2/Images/fig_W20Tg07_OLvsCL_cp.png}
    \caption{Full-order model pressure-coefficient response to a $W_0 = \SI{20}{\meter\per\second}$, $T_g = \SI{0.7}{\second}$ discrete gust (Test~1), with the flap held at trim (open loop) against the same gust with the flap driven by the deflection schedule of the trained Gust Load Alleviation (GLA) controller (closed loop). Pressure coefficient $C_p$ at four instants of the gust; the row beneath each pair is the delta of the pre-gust state, $t = \SI{0}{\second}$, vs that particular $t$. From this sequence it emerges how the flap deflection counteracts the lift increase due to the gust.}
    \label{fig:cfd}
\end{figure}

This is the reason why Gust Load Alleviation (GLA) systems have been developed for several decades and are now part of the design of most modern transport aircraft. Active load alleviation entered service in the 1970s with the C-5A, whose ailerons were driven by accelerometer feedback to relieve the wing-root bending moment~\cite{disney1977c5a}, and later became an enabling technology for load-critical configurations such as the B-2~\cite{britt2000b2}; surveys of the applications of active control to gust alleviation and flutter suppression are given in~\cite{regan2012survey,GLA_review}. The principle is simple: if a control surface, typically a flap or an aileron, is deflected at the right time and with the right amplitude, the aerodynamic load produced by the gust can be partially compensated, and the corresponding peak on the structure can be reduced. During design, a lighter wing can be sized for the reduced load envelope, with a direct impact on weight and fuel consumption; during operation, the same structure can instead be flown safely through gust conditions more severe than its original margins would allow~\cite{regan2012survey}. \Cref{fig:cfd} previews this benefit directly in the full-order model.

GLA architectures differ in the information made available to the controller and in the control law that maps this information onto the control surfaces. Purely reactive feedback schemes act on the structural response measured by accelerometers and inertial sensors, and are inherently limited by the delay between the gust and the measurable response of the structure~\cite{regan2012survey}. Within this class, classical proportional--integral--derivative loops on the measured accelerations have been used to achieve GLA in the wind tunnel, up to the scale of a full elastic aircraft model~\cite{zhao2023windtunnel}, linear quadratic regulators have been applied to gust alleviation on a flexible wing with piezoelectric actuation~\cite{versiani2019smartwing}, and linear quadratic Gaussian designs complement the optimal state feedback with an estimator of the aeroelastic states that cannot be measured directly~\cite{duessler2023lqg}. All these laws are synthesised on a linearised, low-order, low-fidelity model of the aeroelastic system, and their performance degrades when the true response deviates from what that model predicts; robust $H_2$ and $H_\infty$ formulations mitigate this sensitivity by enforcing performance over a prescribed set of model uncertainties, and have been validated in aeroservoelastic wind tunnel tests~\cite{ting2023preview,he2023robust}. Nonlinear and adaptive techniques relax the dependence on an accurate model in a different way: incremental nonlinear dynamic inversion replaces most of the model with acceleration measurements, so that only the control effectiveness needs to be known~\cite{wang2019indi}; dynamic control allocation distributes the alleviation command over redundant control surfaces while accounting for the actuator dynamics~\cite{hansen2022allocation}; active disturbance rejection control lumps the gust and the modelling errors into a total disturbance that an extended state observer estimates and cancels online~\cite{yue2023adrc}.

Feedback alone, however, can only react after the gust is encountered. Feedforward and mixed feedforward--feedback schemes add a direct measurement of the disturbance, taken at the aircraft nose by air-data sensors or ahead of the aircraft by remote sensing~\cite{alam2015mixed}: airborne Doppler lidar can measure the vertical wind field tens of metres ahead of the wing~\cite{rabadan2010lidar}, and lidar-based preview has become the reference architecture for the anticipatory alleviation of discrete gusts~\cite{fezans2017lidar,cavaliere2022estimator,fournier2022lidar}. Model Predictive Control exploits such a preview naturally, because the disturbance forecast enters a receding-horizon optimisation that also enforces the actuator limits explicitly~\cite{rawlings2017mpc}; it has been applied to the GLA of large transport aircraft~\cite{giesseler2012mpc}, also in combination with a linear quadratic Gaussian estimator of the unmeasured states~\cite{liu2017lqgmpc}, and preview has likewise been embedded in robust synthesis~\cite{ting2023preview}. More recently, Machine Learning has been brought into the control law itself, with neural network feedforward compensators trained on gust responses~\cite{zhou2024feedforward} and Reinforcement Learning policies commanding a morphing camber~\cite{hu2025rl}; such controllers capture the nonlinearity of the aeroelastic response, but encode the controlled dynamics implicitly in the policy and give up the explicit constraint handling of a model-based design. With the exception of these model-free schemes, all the architectures above share the same requirement: an internal model of the aeroelastic response that is accurate enough to predict the effect of the control surface and cheap enough to be evaluated thousands of times inside a receding-horizon optimisation.

The design of an active GLA system is based upon two environments. The first one is a model of the aeroelastic response of the wing to the gust, accurate enough to capture the relevant physics, namely the unsteady aerodynamics, the structural dynamics and the coupling between the two. The second one is a control strategy that uses this model to compute, at every instant, how much the flap should deflect in order to minimise a given performance metric. Each of these two ingredients can be addressed at different levels of fidelity, and the two choices are not independent. A full-order model, based on a Computational Fluid Dynamics (CFD) solver coupled with a structural integrator in a Fluid-Structure Interaction (FSI) framework~\cite{bazilevs2013csfi}, captures the physics in great detail but has a computational cost that is incompatible with the repeated evaluation a control law and an iterative design loop both require, and is also unsuitable for control design, because it does not provide a compact representation of the system that the controller can act upon. Low-order models obtained from linear unsteady aerodynamic theories~\cite{Theodorsen,hoblit1988gust}, on the other hand, are cheap enough to sustain that repeated evaluation, but rely on assumptions of small amplitudes, attached flow and inviscid behaviour that are not always justified during a strong gust event, especially when the flap is deflected by a large angle.

In recent years, this gap between the full-order model and control-oriented models has been progressively filled by data-driven surrogate models~\cite{brunton2020ml}. The idea is to run a limited number of full-order simulations, collect the resulting data, and train a surrogate that learns the dynamics of the system from these data. Once trained, the surrogate replaces the original solver in any task for which speed and differentiability matter, including parametric studies, uncertainty quantification and, of particular interest here, model-based control. Many architectures have been proposed for this purpose, ranging from classical projection-based reduced-order models~\cite{BennerGugercinWillcox2015}, such as Proper Orthogonal Decomposition combined with a Galerkin projection~\cite{QuarteroniManzoniNegri2016}, which have been applied to build reduced aeroelastic models for GLA and active flutter suppression~\cite{ZhangChenYe2014,BertolinEtAl2018}, to autoencoder-based models in which a neural network compresses the state of the system into a low-dimensional latent space and a second network propagates it in time~\cite{LeeCarlberg2020,FrescaDedeManzoni2021}, with recent applications to nonlinear flutter prediction on airfoil and wing configurations~\cite{MogharabinEpureanuGhadami2025,WangEtAl2025}. Among the most recent contributions, Latent Dynamics Networks (LDNets)~\cite{Regazzoni}, which extend to spatio-temporal fields the non-intrusive learning of time-dependent dynamics from input-output pairs~\cite{RegazzoniDedeQuarteroni2019}, learn a low-dimensional latent representation and its dynamics directly from input-output trajectories, without the need for a separate autoencoder, and have shown a remarkable accuracy on a range of spatio-temporal problems with a small number of latent variables.

The present work proposes to obtain the internal model of a GLA loop by training it on data produced by the full-order model, instead of deriving it from a simplified aerodynamic theory. The internal model is a LDNet~\cite{Regazzoni}, trained on a campaign of full-order simulations of a two-dimensional airfoil with a hinged trailing-edge flap. The LDNet supplies the aerodynamic loads to a Model Predictive Controller~\cite{rawlings2017mpc}, which is driven by a lidar preview~\cite{rabadan2010lidar,fezans2017lidar} of the incoming gust. The LDNet reproduces the viscous, unsteady and nonlinear response to gust and flap computed by the full-order model, and is used at every control step to find the optimal flap deflection for a receding horizon. It carries neither the cost of the full-order model nor the small-amplitude and attached-flow assumptions of linear unsteady aerodynamic theories. Operating inside the loop constrains how the LDNet is trained. It consumes its own predictions, and it must remain sensitive to the flap deflection commanded by the Model Predictive Controller; neither property follows from a low prediction error measured offline.

This work is organised as follows. In \cref{sec:model} the model is decomposed into subsystems to facilitate its description: the aeroelastic model of the airfoil, the LDNet that surrogates the aerodynamic coefficients and the architecture and cost function of the GLA controller. \Cref{sec:methods} presents the methods: the time discretisation of the closed-loop system, the simulation campaign that generates the training dataset, and the training of the LDNet. \Cref{sec:results} reports the results, from the accuracy of the LDNet to the closed-loop performance over the certification gust envelope and its robustness to measurement noise and systematic errors. \Cref{app:fom} collects the numerical solution of the full-order model, and \cref{app:vv} its verification and validation.
